\documentclass[11pt]{article}

    \usepackage[breakable]{tcolorbox}
    \usepackage{parskip} % Stop auto-indenting (to mimic markdown behaviour)
    

    % Basic figure setup, for now with no caption control since it's done
    % automatically by Pandoc (which extracts ![](path) syntax from Markdown).
    \usepackage{graphicx}
    % Keep aspect ratio if custom image width or height is specified
    \setkeys{Gin}{keepaspectratio}
    % Maintain compatibility with old templates. Remove in nbconvert 6.0
    \let\Oldincludegraphics\includegraphics
    % Ensure that by default, figures have no caption (until we provide a
    % proper Figure object with a Caption API and a way to capture that
    % in the conversion process - todo).
    \usepackage{caption}
    \DeclareCaptionFormat{nocaption}{}
    \captionsetup{format=nocaption,aboveskip=0pt,belowskip=0pt}

    \usepackage{float}
    \floatplacement{figure}{H} % forces figures to be placed at the correct location
    \usepackage{xcolor} % Allow colors to be defined
    \usepackage{enumerate} % Needed for markdown enumerations to work
    \usepackage{geometry} % Used to adjust the document margins
    \usepackage{amsmath} % Equations
    \usepackage{amssymb} % Equations
    \usepackage{textcomp} % defines textquotesingle
    % Hack from http://tex.stackexchange.com/a/47451/13684:
    \AtBeginDocument{%
        \def\PYZsq{\textquotesingle}% Upright quotes in Pygmentized code
    }
    \usepackage{upquote} % Upright quotes for verbatim code
    \usepackage{eurosym} % defines \euro

    \usepackage{iftex}
    \ifPDFTeX
        \usepackage[T1]{fontenc}
        \IfFileExists{alphabeta.sty}{
              \usepackage{alphabeta}
          }{
              \usepackage[mathletters]{ucs}
              \usepackage[utf8x]{inputenc}
          }
    \else
        \usepackage{fontspec}
        \usepackage{unicode-math}
    \fi

    \usepackage{fancyvrb} % verbatim replacement that allows latex
    \usepackage{grffile} % extends the file name processing of package graphics
                         % to support a larger range
    \makeatletter % fix for old versions of grffile with XeLaTeX
    \@ifpackagelater{grffile}{2019/11/01}
    {
      % Do nothing on new versions
    }
    {
      \def\Gread@@xetex#1{%
        \IfFileExists{"\Gin@base".bb}%
        {\Gread@eps{\Gin@base.bb}}%
        {\Gread@@xetex@aux#1}%
      }
    }
    \makeatother
    \usepackage[Export]{adjustbox} % Used to constrain images to a maximum size
    \adjustboxset{max size={0.9\linewidth}{0.9\paperheight}}

    % The hyperref package gives us a pdf with properly built
    % internal navigation ('pdf bookmarks' for the table of contents,
    % internal cross-reference links, web links for URLs, etc.)
    \usepackage{hyperref}
    % The default LaTeX title has an obnoxious amount of whitespace. By default,
    % titling removes some of it. It also provides customization options.
    \usepackage{titling}
    \usepackage{longtable} % longtable support required by pandoc >1.10
    \usepackage{booktabs}  % table support for pandoc > 1.12.2
    \usepackage{array}     % table support for pandoc >= 2.11.3
    \usepackage{calc}      % table minipage width calculation for pandoc >= 2.11.1
    \usepackage[inline]{enumitem} % IRkernel/repr support (it uses the enumerate* environment)
    \usepackage[normalem]{ulem} % ulem is needed to support strikethroughs (\sout)
                                % normalem makes italics be italics, not underlines
    \usepackage{soul}      % strikethrough (\st) support for pandoc >= 3.0.0
    \usepackage{mathrsfs}
    

    
    % Colors for the hyperref package
    \definecolor{urlcolor}{rgb}{0,.145,.698}
    \definecolor{linkcolor}{rgb}{.71,0.21,0.01}
    \definecolor{citecolor}{rgb}{.12,.54,.11}

    % ANSI colors
    \definecolor{ansi-black}{HTML}{3E424D}
    \definecolor{ansi-black-intense}{HTML}{282C36}
    \definecolor{ansi-red}{HTML}{E75C58}
    \definecolor{ansi-red-intense}{HTML}{B22B31}
    \definecolor{ansi-green}{HTML}{00A250}
    \definecolor{ansi-green-intense}{HTML}{007427}
    \definecolor{ansi-yellow}{HTML}{DDB62B}
    \definecolor{ansi-yellow-intense}{HTML}{B27D12}
    \definecolor{ansi-blue}{HTML}{208FFB}
    \definecolor{ansi-blue-intense}{HTML}{0065CA}
    \definecolor{ansi-magenta}{HTML}{D160C4}
    \definecolor{ansi-magenta-intense}{HTML}{A03196}
    \definecolor{ansi-cyan}{HTML}{60C6C8}
    \definecolor{ansi-cyan-intense}{HTML}{258F8F}
    \definecolor{ansi-white}{HTML}{C5C1B4}
    \definecolor{ansi-white-intense}{HTML}{A1A6B2}
    \definecolor{ansi-default-inverse-fg}{HTML}{FFFFFF}
    \definecolor{ansi-default-inverse-bg}{HTML}{000000}

    % common color for the border for error outputs.
    \definecolor{outerrorbackground}{HTML}{FFDFDF}

    % commands and environments needed by pandoc snippets
    % extracted from the output of `pandoc -s`
    \providecommand{\tightlist}{%
      \setlength{\itemsep}{0pt}\setlength{\parskip}{0pt}}
    \DefineVerbatimEnvironment{Highlighting}{Verbatim}{commandchars=\\\{\}}
    % Add ',fontsize=\small' for more characters per line
    \newenvironment{Shaded}{}{}
    \newcommand{\KeywordTok}[1]{\textcolor[rgb]{0.00,0.44,0.13}{\textbf{{#1}}}}
    \newcommand{\DataTypeTok}[1]{\textcolor[rgb]{0.56,0.13,0.00}{{#1}}}
    \newcommand{\DecValTok}[1]{\textcolor[rgb]{0.25,0.63,0.44}{{#1}}}
    \newcommand{\BaseNTok}[1]{\textcolor[rgb]{0.25,0.63,0.44}{{#1}}}
    \newcommand{\FloatTok}[1]{\textcolor[rgb]{0.25,0.63,0.44}{{#1}}}
    \newcommand{\CharTok}[1]{\textcolor[rgb]{0.25,0.44,0.63}{{#1}}}
    \newcommand{\StringTok}[1]{\textcolor[rgb]{0.25,0.44,0.63}{{#1}}}
    \newcommand{\CommentTok}[1]{\textcolor[rgb]{0.38,0.63,0.69}{\textit{{#1}}}}
    \newcommand{\OtherTok}[1]{\textcolor[rgb]{0.00,0.44,0.13}{{#1}}}
    \newcommand{\AlertTok}[1]{\textcolor[rgb]{1.00,0.00,0.00}{\textbf{{#1}}}}
    \newcommand{\FunctionTok}[1]{\textcolor[rgb]{0.02,0.16,0.49}{{#1}}}
    \newcommand{\RegionMarkerTok}[1]{{#1}}
    \newcommand{\ErrorTok}[1]{\textcolor[rgb]{1.00,0.00,0.00}{\textbf{{#1}}}}
    \newcommand{\NormalTok}[1]{{#1}}

    % Additional commands for more recent versions of Pandoc
    \newcommand{\ConstantTok}[1]{\textcolor[rgb]{0.53,0.00,0.00}{{#1}}}
    \newcommand{\SpecialCharTok}[1]{\textcolor[rgb]{0.25,0.44,0.63}{{#1}}}
    \newcommand{\VerbatimStringTok}[1]{\textcolor[rgb]{0.25,0.44,0.63}{{#1}}}
    \newcommand{\SpecialStringTok}[1]{\textcolor[rgb]{0.73,0.40,0.53}{{#1}}}
    \newcommand{\ImportTok}[1]{{#1}}
    \newcommand{\DocumentationTok}[1]{\textcolor[rgb]{0.73,0.13,0.13}{\textit{{#1}}}}
    \newcommand{\AnnotationTok}[1]{\textcolor[rgb]{0.38,0.63,0.69}{\textbf{\textit{{#1}}}}}
    \newcommand{\CommentVarTok}[1]{\textcolor[rgb]{0.38,0.63,0.69}{\textbf{\textit{{#1}}}}}
    \newcommand{\VariableTok}[1]{\textcolor[rgb]{0.10,0.09,0.49}{{#1}}}
    \newcommand{\ControlFlowTok}[1]{\textcolor[rgb]{0.00,0.44,0.13}{\textbf{{#1}}}}
    \newcommand{\OperatorTok}[1]{\textcolor[rgb]{0.40,0.40,0.40}{{#1}}}
    \newcommand{\BuiltInTok}[1]{{#1}}
    \newcommand{\ExtensionTok}[1]{{#1}}
    \newcommand{\PreprocessorTok}[1]{\textcolor[rgb]{0.74,0.48,0.00}{{#1}}}
    \newcommand{\AttributeTok}[1]{\textcolor[rgb]{0.49,0.56,0.16}{{#1}}}
    \newcommand{\InformationTok}[1]{\textcolor[rgb]{0.38,0.63,0.69}{\textbf{\textit{{#1}}}}}
    \newcommand{\WarningTok}[1]{\textcolor[rgb]{0.38,0.63,0.69}{\textbf{\textit{{#1}}}}}
    \makeatletter
    \newsavebox\pandoc@box
    \newcommand*\pandocbounded[1]{%
      \sbox\pandoc@box{#1}%
      % scaling factors for width and height
      \Gscale@div\@tempa\textheight{\dimexpr\ht\pandoc@box+\dp\pandoc@box\relax}%
      \Gscale@div\@tempb\linewidth{\wd\pandoc@box}%
      % select the smaller of both
      \ifdim\@tempb\p@<\@tempa\p@
        \let\@tempa\@tempb
      \fi
      % scaling accordingly (\@tempa < 1)
      \ifdim\@tempa\p@<\p@
        \scalebox{\@tempa}{\usebox\pandoc@box}%
      % scaling not needed, use as it is
      \else
        \usebox{\pandoc@box}%
      \fi
    }
    \makeatother

    % Define a nice break command that doesn't care if a line doesn't already
    % exist.
    \def\br{\hspace*{\fill} \\* }
    % Math Jax compatibility definitions
    \def\gt{>}
    \def\lt{<}
    \let\Oldtex\TeX
    \let\Oldlatex\LaTeX
    \renewcommand{\TeX}{\textrm{\Oldtex}}
    \renewcommand{\LaTeX}{\textrm{\Oldlatex}}
    % Document parameters
    % Document title
    \title{587-N13b}
    
    
    
    
    
    
    
% Pygments definitions
\makeatletter
\def\PY@reset{\let\PY@it=\relax \let\PY@bf=\relax%
    \let\PY@ul=\relax \let\PY@tc=\relax%
    \let\PY@bc=\relax \let\PY@ff=\relax}
\def\PY@tok#1{\csname PY@tok@#1\endcsname}
\def\PY@toks#1+{\ifx\relax#1\empty\else%
    \PY@tok{#1}\expandafter\PY@toks\fi}
\def\PY@do#1{\PY@bc{\PY@tc{\PY@ul{%
    \PY@it{\PY@bf{\PY@ff{#1}}}}}}}
\def\PY#1#2{\PY@reset\PY@toks#1+\relax+\PY@do{#2}}

\@namedef{PY@tok@w}{\def\PY@tc##1{\textcolor[rgb]{0.73,0.73,0.73}{##1}}}
\@namedef{PY@tok@c}{\let\PY@it=\textit\def\PY@tc##1{\textcolor[rgb]{0.24,0.48,0.48}{##1}}}
\@namedef{PY@tok@cp}{\def\PY@tc##1{\textcolor[rgb]{0.61,0.40,0.00}{##1}}}
\@namedef{PY@tok@k}{\let\PY@bf=\textbf\def\PY@tc##1{\textcolor[rgb]{0.00,0.50,0.00}{##1}}}
\@namedef{PY@tok@kp}{\def\PY@tc##1{\textcolor[rgb]{0.00,0.50,0.00}{##1}}}
\@namedef{PY@tok@kt}{\def\PY@tc##1{\textcolor[rgb]{0.69,0.00,0.25}{##1}}}
\@namedef{PY@tok@o}{\def\PY@tc##1{\textcolor[rgb]{0.40,0.40,0.40}{##1}}}
\@namedef{PY@tok@ow}{\let\PY@bf=\textbf\def\PY@tc##1{\textcolor[rgb]{0.67,0.13,1.00}{##1}}}
\@namedef{PY@tok@nb}{\def\PY@tc##1{\textcolor[rgb]{0.00,0.50,0.00}{##1}}}
\@namedef{PY@tok@nf}{\def\PY@tc##1{\textcolor[rgb]{0.00,0.00,1.00}{##1}}}
\@namedef{PY@tok@nc}{\let\PY@bf=\textbf\def\PY@tc##1{\textcolor[rgb]{0.00,0.00,1.00}{##1}}}
\@namedef{PY@tok@nn}{\let\PY@bf=\textbf\def\PY@tc##1{\textcolor[rgb]{0.00,0.00,1.00}{##1}}}
\@namedef{PY@tok@ne}{\let\PY@bf=\textbf\def\PY@tc##1{\textcolor[rgb]{0.80,0.25,0.22}{##1}}}
\@namedef{PY@tok@nv}{\def\PY@tc##1{\textcolor[rgb]{0.10,0.09,0.49}{##1}}}
\@namedef{PY@tok@no}{\def\PY@tc##1{\textcolor[rgb]{0.53,0.00,0.00}{##1}}}
\@namedef{PY@tok@nl}{\def\PY@tc##1{\textcolor[rgb]{0.46,0.46,0.00}{##1}}}
\@namedef{PY@tok@ni}{\let\PY@bf=\textbf\def\PY@tc##1{\textcolor[rgb]{0.44,0.44,0.44}{##1}}}
\@namedef{PY@tok@na}{\def\PY@tc##1{\textcolor[rgb]{0.41,0.47,0.13}{##1}}}
\@namedef{PY@tok@nt}{\let\PY@bf=\textbf\def\PY@tc##1{\textcolor[rgb]{0.00,0.50,0.00}{##1}}}
\@namedef{PY@tok@nd}{\def\PY@tc##1{\textcolor[rgb]{0.67,0.13,1.00}{##1}}}
\@namedef{PY@tok@s}{\def\PY@tc##1{\textcolor[rgb]{0.73,0.13,0.13}{##1}}}
\@namedef{PY@tok@sd}{\let\PY@it=\textit\def\PY@tc##1{\textcolor[rgb]{0.73,0.13,0.13}{##1}}}
\@namedef{PY@tok@si}{\let\PY@bf=\textbf\def\PY@tc##1{\textcolor[rgb]{0.64,0.35,0.47}{##1}}}
\@namedef{PY@tok@se}{\let\PY@bf=\textbf\def\PY@tc##1{\textcolor[rgb]{0.67,0.36,0.12}{##1}}}
\@namedef{PY@tok@sr}{\def\PY@tc##1{\textcolor[rgb]{0.64,0.35,0.47}{##1}}}
\@namedef{PY@tok@ss}{\def\PY@tc##1{\textcolor[rgb]{0.10,0.09,0.49}{##1}}}
\@namedef{PY@tok@sx}{\def\PY@tc##1{\textcolor[rgb]{0.00,0.50,0.00}{##1}}}
\@namedef{PY@tok@m}{\def\PY@tc##1{\textcolor[rgb]{0.40,0.40,0.40}{##1}}}
\@namedef{PY@tok@gh}{\let\PY@bf=\textbf\def\PY@tc##1{\textcolor[rgb]{0.00,0.00,0.50}{##1}}}
\@namedef{PY@tok@gu}{\let\PY@bf=\textbf\def\PY@tc##1{\textcolor[rgb]{0.50,0.00,0.50}{##1}}}
\@namedef{PY@tok@gd}{\def\PY@tc##1{\textcolor[rgb]{0.63,0.00,0.00}{##1}}}
\@namedef{PY@tok@gi}{\def\PY@tc##1{\textcolor[rgb]{0.00,0.52,0.00}{##1}}}
\@namedef{PY@tok@gr}{\def\PY@tc##1{\textcolor[rgb]{0.89,0.00,0.00}{##1}}}
\@namedef{PY@tok@ge}{\let\PY@it=\textit}
\@namedef{PY@tok@gs}{\let\PY@bf=\textbf}
\@namedef{PY@tok@ges}{\let\PY@bf=\textbf\let\PY@it=\textit}
\@namedef{PY@tok@gp}{\let\PY@bf=\textbf\def\PY@tc##1{\textcolor[rgb]{0.00,0.00,0.50}{##1}}}
\@namedef{PY@tok@go}{\def\PY@tc##1{\textcolor[rgb]{0.44,0.44,0.44}{##1}}}
\@namedef{PY@tok@gt}{\def\PY@tc##1{\textcolor[rgb]{0.00,0.27,0.87}{##1}}}
\@namedef{PY@tok@err}{\def\PY@bc##1{{\setlength{\fboxsep}{\string -\fboxrule}\fcolorbox[rgb]{1.00,0.00,0.00}{1,1,1}{\strut ##1}}}}
\@namedef{PY@tok@kc}{\let\PY@bf=\textbf\def\PY@tc##1{\textcolor[rgb]{0.00,0.50,0.00}{##1}}}
\@namedef{PY@tok@kd}{\let\PY@bf=\textbf\def\PY@tc##1{\textcolor[rgb]{0.00,0.50,0.00}{##1}}}
\@namedef{PY@tok@kn}{\let\PY@bf=\textbf\def\PY@tc##1{\textcolor[rgb]{0.00,0.50,0.00}{##1}}}
\@namedef{PY@tok@kr}{\let\PY@bf=\textbf\def\PY@tc##1{\textcolor[rgb]{0.00,0.50,0.00}{##1}}}
\@namedef{PY@tok@bp}{\def\PY@tc##1{\textcolor[rgb]{0.00,0.50,0.00}{##1}}}
\@namedef{PY@tok@fm}{\def\PY@tc##1{\textcolor[rgb]{0.00,0.00,1.00}{##1}}}
\@namedef{PY@tok@vc}{\def\PY@tc##1{\textcolor[rgb]{0.10,0.09,0.49}{##1}}}
\@namedef{PY@tok@vg}{\def\PY@tc##1{\textcolor[rgb]{0.10,0.09,0.49}{##1}}}
\@namedef{PY@tok@vi}{\def\PY@tc##1{\textcolor[rgb]{0.10,0.09,0.49}{##1}}}
\@namedef{PY@tok@vm}{\def\PY@tc##1{\textcolor[rgb]{0.10,0.09,0.49}{##1}}}
\@namedef{PY@tok@sa}{\def\PY@tc##1{\textcolor[rgb]{0.73,0.13,0.13}{##1}}}
\@namedef{PY@tok@sb}{\def\PY@tc##1{\textcolor[rgb]{0.73,0.13,0.13}{##1}}}
\@namedef{PY@tok@sc}{\def\PY@tc##1{\textcolor[rgb]{0.73,0.13,0.13}{##1}}}
\@namedef{PY@tok@dl}{\def\PY@tc##1{\textcolor[rgb]{0.73,0.13,0.13}{##1}}}
\@namedef{PY@tok@s2}{\def\PY@tc##1{\textcolor[rgb]{0.73,0.13,0.13}{##1}}}
\@namedef{PY@tok@sh}{\def\PY@tc##1{\textcolor[rgb]{0.73,0.13,0.13}{##1}}}
\@namedef{PY@tok@s1}{\def\PY@tc##1{\textcolor[rgb]{0.73,0.13,0.13}{##1}}}
\@namedef{PY@tok@mb}{\def\PY@tc##1{\textcolor[rgb]{0.40,0.40,0.40}{##1}}}
\@namedef{PY@tok@mf}{\def\PY@tc##1{\textcolor[rgb]{0.40,0.40,0.40}{##1}}}
\@namedef{PY@tok@mh}{\def\PY@tc##1{\textcolor[rgb]{0.40,0.40,0.40}{##1}}}
\@namedef{PY@tok@mi}{\def\PY@tc##1{\textcolor[rgb]{0.40,0.40,0.40}{##1}}}
\@namedef{PY@tok@il}{\def\PY@tc##1{\textcolor[rgb]{0.40,0.40,0.40}{##1}}}
\@namedef{PY@tok@mo}{\def\PY@tc##1{\textcolor[rgb]{0.40,0.40,0.40}{##1}}}
\@namedef{PY@tok@ch}{\let\PY@it=\textit\def\PY@tc##1{\textcolor[rgb]{0.24,0.48,0.48}{##1}}}
\@namedef{PY@tok@cm}{\let\PY@it=\textit\def\PY@tc##1{\textcolor[rgb]{0.24,0.48,0.48}{##1}}}
\@namedef{PY@tok@cpf}{\let\PY@it=\textit\def\PY@tc##1{\textcolor[rgb]{0.24,0.48,0.48}{##1}}}
\@namedef{PY@tok@c1}{\let\PY@it=\textit\def\PY@tc##1{\textcolor[rgb]{0.24,0.48,0.48}{##1}}}
\@namedef{PY@tok@cs}{\let\PY@it=\textit\def\PY@tc##1{\textcolor[rgb]{0.24,0.48,0.48}{##1}}}

\def\PYZbs{\char`\\}
\def\PYZus{\char`\_}
\def\PYZob{\char`\{}
\def\PYZcb{\char`\}}
\def\PYZca{\char`\^}
\def\PYZam{\char`\&}
\def\PYZlt{\char`\<}
\def\PYZgt{\char`\>}
\def\PYZsh{\char`\#}
\def\PYZpc{\char`\%}
\def\PYZdl{\char`\$}
\def\PYZhy{\char`\-}
\def\PYZsq{\char`\'}
\def\PYZdq{\char`\"}
\def\PYZti{\char`\~}
% for compatibility with earlier versions
\def\PYZat{@}
\def\PYZlb{[}
\def\PYZrb{]}
\makeatother


    % For linebreaks inside Verbatim environment from package fancyvrb.
    \makeatletter
        \newbox\Wrappedcontinuationbox
        \newbox\Wrappedvisiblespacebox
        \newcommand*\Wrappedvisiblespace {\textcolor{red}{\textvisiblespace}}
        \newcommand*\Wrappedcontinuationsymbol {\textcolor{red}{\llap{\tiny$\m@th\hookrightarrow$}}}
        \newcommand*\Wrappedcontinuationindent {3ex }
        \newcommand*\Wrappedafterbreak {\kern\Wrappedcontinuationindent\copy\Wrappedcontinuationbox}
        % Take advantage of the already applied Pygments mark-up to insert
        % potential linebreaks for TeX processing.
        %        {, <, #, %, $, ' and ": go to next line.
        %        _, }, ^, &, >, - and ~: stay at end of broken line.
        % Use of \textquotesingle for straight quote.
        \newcommand*\Wrappedbreaksatspecials {%
            \def\PYGZus{\discretionary{\char`\_}{\Wrappedafterbreak}{\char`\_}}%
            \def\PYGZob{\discretionary{}{\Wrappedafterbreak\char`\{}{\char`\{}}%
            \def\PYGZcb{\discretionary{\char`\}}{\Wrappedafterbreak}{\char`\}}}%
            \def\PYGZca{\discretionary{\char`\^}{\Wrappedafterbreak}{\char`\^}}%
            \def\PYGZam{\discretionary{\char`\&}{\Wrappedafterbreak}{\char`\&}}%
            \def\PYGZlt{\discretionary{}{\Wrappedafterbreak\char`\<}{\char`\<}}%
            \def\PYGZgt{\discretionary{\char`\>}{\Wrappedafterbreak}{\char`\>}}%
            \def\PYGZsh{\discretionary{}{\Wrappedafterbreak\char`\#}{\char`\#}}%
            \def\PYGZpc{\discretionary{}{\Wrappedafterbreak\char`\%}{\char`\%}}%
            \def\PYGZdl{\discretionary{}{\Wrappedafterbreak\char`\$}{\char`\$}}%
            \def\PYGZhy{\discretionary{\char`\-}{\Wrappedafterbreak}{\char`\-}}%
            \def\PYGZsq{\discretionary{}{\Wrappedafterbreak\textquotesingle}{\textquotesingle}}%
            \def\PYGZdq{\discretionary{}{\Wrappedafterbreak\char`\"}{\char`\"}}%
            \def\PYGZti{\discretionary{\char`\~}{\Wrappedafterbreak}{\char`\~}}%
        }
        % Some characters . , ; ? ! / are not pygmentized.
        % This macro makes them "active" and they will insert potential linebreaks
        \newcommand*\Wrappedbreaksatpunct {%
            \lccode`\~`\.\lowercase{\def~}{\discretionary{\hbox{\char`\.}}{\Wrappedafterbreak}{\hbox{\char`\.}}}%
            \lccode`\~`\,\lowercase{\def~}{\discretionary{\hbox{\char`\,}}{\Wrappedafterbreak}{\hbox{\char`\,}}}%
            \lccode`\~`\;\lowercase{\def~}{\discretionary{\hbox{\char`\;}}{\Wrappedafterbreak}{\hbox{\char`\;}}}%
            \lccode`\~`\:\lowercase{\def~}{\discretionary{\hbox{\char`\:}}{\Wrappedafterbreak}{\hbox{\char`\:}}}%
            \lccode`\~`\?\lowercase{\def~}{\discretionary{\hbox{\char`\?}}{\Wrappedafterbreak}{\hbox{\char`\?}}}%
            \lccode`\~`\!\lowercase{\def~}{\discretionary{\hbox{\char`\!}}{\Wrappedafterbreak}{\hbox{\char`\!}}}%
            \lccode`\~`\/\lowercase{\def~}{\discretionary{\hbox{\char`\/}}{\Wrappedafterbreak}{\hbox{\char`\/}}}%
            \catcode`\.\active
            \catcode`\,\active
            \catcode`\;\active
            \catcode`\:\active
            \catcode`\?\active
            \catcode`\!\active
            \catcode`\/\active
            \lccode`\~`\~
        }
    \makeatother

    \let\OriginalVerbatim=\Verbatim
    \makeatletter
    \renewcommand{\Verbatim}[1][1]{%
        %\parskip\z@skip
        \sbox\Wrappedcontinuationbox {\Wrappedcontinuationsymbol}%
        \sbox\Wrappedvisiblespacebox {\FV@SetupFont\Wrappedvisiblespace}%
        \def\FancyVerbFormatLine ##1{\hsize\linewidth
            \vtop{\raggedright\hyphenpenalty\z@\exhyphenpenalty\z@
                \doublehyphendemerits\z@\finalhyphendemerits\z@
                \strut ##1\strut}%
        }%
        % If the linebreak is at a space, the latter will be displayed as visible
        % space at end of first line, and a continuation symbol starts next line.
        % Stretch/shrink are however usually zero for typewriter font.
        \def\FV@Space {%
            \nobreak\hskip\z@ plus\fontdimen3\font minus\fontdimen4\font
            \discretionary{\copy\Wrappedvisiblespacebox}{\Wrappedafterbreak}
            {\kern\fontdimen2\font}%
        }%

        % Allow breaks at special characters using \PYG... macros.
        \Wrappedbreaksatspecials
        % Breaks at punctuation characters . , ; ? ! and / need catcode=\active
        \OriginalVerbatim[#1,codes*=\Wrappedbreaksatpunct]%
    }
    \makeatother

    % Exact colors from NB
    \definecolor{incolor}{HTML}{303F9F}
    \definecolor{outcolor}{HTML}{D84315}
    \definecolor{cellborder}{HTML}{CFCFCF}
    \definecolor{cellbackground}{HTML}{F7F7F7}

    % prompt
    \makeatletter
    \newcommand{\boxspacing}{\kern\kvtcb@left@rule\kern\kvtcb@boxsep}
    \makeatother
    \newcommand{\prompt}[4]{
        {\ttfamily\llap{{\color{#2}[#3]:\hspace{3pt}#4}}\vspace{-\baselineskip}}
    }
    

    
    % Prevent overflowing lines due to hard-to-break entities
    \sloppy
    % Setup hyperref package
    \hypersetup{
      breaklinks=true,  % so long urls are correctly broken across lines
      colorlinks=true,
      urlcolor=urlcolor,
      linkcolor=linkcolor,
      citecolor=citecolor,
      }
    % Slightly bigger margins than the latex defaults
    
    \geometry{verbose,tmargin=1in,bmargin=1in,lmargin=1in,rmargin=1in}
    
    

\begin{document}
    
    \maketitle
    
    

    
    \hypertarget{material-balances-i}{%
\section{Material Balances I}\label{material-balances-i}}

\href{https://www.youtube.com/watch?v=7eXTF6paB6Q\&list=RD7eXTF6paB6Q\&start_radio=1}{What
is Bond listening to right now?}

This Notebook contrasts intensive and extensive properties of systems,
and it introduces material balances on reacting systems.

    \hypertarget{intensive-and-extensive-properties}{%
\subsection{Intensive and Extensive
Properties}\label{intensive-and-extensive-properties}}

It is worth a brief discussion of intensive properties and extensive
properties. We have to relate the two frequently in reactor design, so
it is important that we have good intuition about what they are and how
they are different. Below, we'll establish a working conceptual
definition of these two things:

\hypertarget{intensive-properties}{%
\subsubsection{Intensive Properties}\label{intensive-properties}}

\textbf{Intensive properties} are scale-independent. They are presented
as values per unit of something (e.g., per unit volume, per unit mass,
per unit area, etc.). Some examples would be concentration (moles per
volume); density (mass per volume); mole fraction (moles per total
moles); etc. We also typically report reaction rates as intensive
quantities that are normalized to volume, catalyst mass, or to some
other characteristic dimension of the system. Intensive properties are
the same whether you're describing a large reactor or a small reactor.
For example, a 1 Molar solution of sulfuric acid always has a
concentration of 1 mole \(H_2SO_4\) per liter of solution volume whether
you are considering that sulfuric acid in a 1L system or a 100L system.
Intensive properties are extremely useful because they are system/scale
independent.

\hypertarget{extensive-properties}{%
\subsubsection{Extensive Properties}\label{extensive-properties}}

\textbf{Extensive Properties} are scale-dependent. They will not be
presented as being normalized to a unit value of something. Some
examples include total moles in a reactor, molar flowrates of species
entering and leaving a reactor, volumes of reactors, and the mass of
catalyst inside of a reactor. These quantities change as you consider
processes of different sizes or scales. They are critically important
because they define the actual physical size of our system, throughput,
revenue, cost, and profitability.

A good test to determine whether a property is intensive or extensive is
to ask ``If I double the size of the system, or if I cut the size of the
system in half, does the value of this property change?'' If it does not
change, it is an intensive property. If it does change, it is an
extensive property.

\hypertarget{the-relationship-between-intensive-and-extensive-properties}{%
\subsubsection{The Relationship Between Intensive and Extensive
Properties}\label{the-relationship-between-intensive-and-extensive-properties}}

For an illustration of intensive and extensive properties, consider a 1L
tank filled with a 1 molar solution of species \(A\), i.e., \(C_A = 1\)
mole/liter. If we wanted to know the total moles of \(A\) in that tank,
we would calculate it as:

\[N_A = C_AV\]

So for this case, we would have 1 mole of \(A\) in our 1L reactor. If we
take the same solution and fill a 100L reactor with it, we would find
that there are 100 total moles of \(A\) in the 100L reactor. The
concentration of \(A\), \(C_A\), is the same in both cases (1M).
Concentration is an intensive variable, and it is independent of scale
(the size of the system). The total number of moles of \(A\), \(N_A\),
is an extensive variable. As we fill larger and larger tanks with a 1M
solution of \(A\), we will get larger and larger quantities of moles of
\(A\). If you get this basic concept, this is an adequate working
understanding of intensive and extensive properties and how both might
show up in a reactor design problem.

    \hypertarget{how-do-reactors-scale-anyway}{%
\subsection{How do reactors scale
anyway?}\label{how-do-reactors-scale-anyway}}

Our courses have primarily conditioned us to think of intensive
properties in units per volume, i.e., they are properties specified per
unit volume of a mixture. We also usually think about the size of the
reactor as a volume, but it is important to consider that the
\textbf{\emph{extensive productivity}} of a reactor (e.g., how many
moles it produces annually) does not necessarily scale closely with the
volume of the reactor. Sometimes, it will scale more strongly with the
mass of catalyst in the reactor, the interfacial area between two
separate phases, or the number of accessible catalytic active sites in
our system. To understand the critical ``size'' dimension of a reactor,
we should first consider the differences between homogeneous reactions
and heterogeneous reactions.

In a \textbf{Homogeneous Reaction} all species--reactants,
intermediates, products, and/or catalysts--are present in a single
phase. Usually, this will be either a gas or a solution phase. Note that
when we say solution, we mean a single liquid phase wherein all species
are miscible and/or fully soluble.

In a \textbf{Heterogeneous Reaction} at least one species--reactant,
intermediate, product, or catalyst--is present in a second phase. Note
that two seperate liquid phases can be an example of a heterogeneous
reaction.

    \hypertarget{homogeneous-reactions}{%
\subsubsection{Homogeneous Reactions}\label{homogeneous-reactions}}

One example of a homogeneous reaction is ethane cracking, which is the
dominant process for converting ethane from oil and gas wells into the
far more valuable commodity chemical, ethylene:

\[C_2H_6 \ (g) \leftrightharpoons C_2H_4 \ (g) + H_2 \ (g)\]

This reaction is non-catalytic. It occurs in the gas phase. The
mechanism of this reaction involves the formation of radical
intermediates, which are also present in the gas phase. Another example
of a liquid-phase homogeneous reaction is glucose isomerization:

\[C_6H_{12}O_6 \ (\textrm{glucose}) \ (aq.) \leftrightharpoons C_6H_{12}O_6 \ (\textrm{fructose}) \ (aq.)\]

This reaction occurs in aqueous solution (i.e., in water), and it uses a
water-soluble enzyme (xylose isomerase) as a catalyst. All of the
intermediates that form in the glucose isomerization mechanism are in
aqueous solution, so we classify this as a homogeneous reaction. An
important thing to recognize about a homogeneous reaction is that it
occurs \textbf{\emph{within the fluid}}. As a consequence, it is
happening throughout the space (volume) occupied by the fluid. One can
then imagine that the \textbf{\emph{extensive}} quantities of species
produced in a homogeneous reaction will scale with the
\textbf{\emph{volume}} of fluid that is reacting.

\hypertarget{a-thought-experiment}{%
\paragraph{A thought experiment}\label{a-thought-experiment}}

As an example, let's say we study ethane pyrolysis to produce ethylene
\((C_2H_6 \ (g) \leftrightharpoons C_2H_4 \ (g) + H_2 \ (g))\) in a
small scale batch reactor that has a volume of 1L. Under these
conditions, the rate of reaction is constant at 1 mole of extent per
liter per minute (this will not normally be true, but we will assume it
here to keep the thought experiment simple). We fill that reactor with 1
mole of ethane so that it has a concentration of 1M. The ethane will
expand to fill the entire reactor volume, and the reaction occurs
throughout the volume occupied. If we assume the reaction is
irreversible, and we let this system react for one minute, we would make
1 mole of ethylene \((C_2H_4)\) and 1 mole of \(H_2\)

Now, let's consider this same reaction in a 1000L reactor, and we run it
under the same temperature, pressure, and intensive condition (i.e., an
ethane concentration of 1M). To do this, we would initially charge this
1000L reactor with an extensive quantity of 1000 moles of ethane
\((C_2H_6)\) so that the ethane concentration is 1M. The reaction again
occurs throughout the entire volume of the reactor, and we again have
the same, constant intensive reaction rate of 1 mole of extent per liter
per minute. If we let this system react for 1 minute in this 1000L
reactor, we would form 1000 moles of each product. Contrast this to the
case of the 1L reactor, which produces only 1 mole of product under
identical conditions. In this case, as we increase the volume of the
reactor, the reaction media fills that volume. The extensive
productivity of the reactor (total moles produced) increases
accordingly, whereas intensive properties (rate per volume,
concentration of ethane) are identical in the two systems because they
are scale independent.

\begin{figure}
\centering
\includegraphics{Homogeneous.png}
\caption{Conceptual Illustration of a Homogeneous Reaction}
\end{figure}

Because \textbf{\emph{homogeneous}} reactions tend to occur throughout
the space occupied by the fluid phase, reason, when we consider
\textbf{\emph{homogeneous}} reactions, we tend to think of the reactor
size in terms of the volume occupied by the fluid phase.

\hypertarget{heterogeneous-reactions}{%
\subsubsection{Heterogeneous Reactions}\label{heterogeneous-reactions}}

Heterogeneous reactions are much more varied. We can imagine lots of
cases where a second phase of matter is present. A classic one is a
gas-phase or liquid-phase reaction catalyzed by a solid catalyst. We can
also imagine a reaction between two immiscible liquids (See Equilibrium
Examples from
\href{https://jqbond.github.io/CEN587/Notebooks/587-N12.html}{Notebook
12}), or a chemical vapor deposition process, wherein gas-phase species
react to form solid products. We'll go through a few examples below.

\hypertarget{catalytic-reactions-at-surfaces-ammonia-synthesis}{%
\paragraph{Catalytic Reactions at Surfaces: Ammonia
Synthesis}\label{catalytic-reactions-at-surfaces-ammonia-synthesis}}

Superficially, this looks like a gas-phase reaction, but it is actually
catalyzed by (and occuring on) a solid surface. So, here, we have gas
phase reactants (\(N_2\) and \(H_2\)) and a gas phase product
(\(N\!H_3\)), but we have a solid catalyst (usually an Iron-based
material), and we have intermediates that are chemically adsorbed onto
the atoms on that metal surface. So, in this case, the catalyst and the
intermediates are all present as a solid phase, and we have a
heterogeneous reaction.

\begin{figure}
\centering
\includegraphics{Ammonia.png}
\caption{Ammonia Synthesis: A Heterogeneous Reaction}
\end{figure}

Here, the reaction is \textbf{\emph{only}} occurring on the catalyst. It
does not occur throughout the space occupied by the fluid, i.e., the
reactor volume. In this case, the critical component of the reactor that
dictates its extensive productivity is the quantity of catalyst inside
of the reactor, not the volume of the reactor itself. Often, it is
convenient to quantify amount of catalyst in terms of its mass;
therefore, mass of catalysis is a common scaling dimension for reactions
that occur on solid surfaces. For an illustration, consider the thought
experiment below.

    \hypertarget{scaling-a-heterogeneously-catalyzed-reaction}{%
\paragraph{Scaling a Heterogeneously catalyzed
reaction}\label{scaling-a-heterogeneously-catalyzed-reaction}}

This is a relatively loose conceptual illustration because it neglects
the dependence of reaction rate on species concentration, but it gets to
the point about how the definition of ``size'' of a system really
depends on the type of reaction we're considering. In reality, changing
the volume of a catalytic reactor will usually have an impact on the
conversion achieved in that reactor, it just won't be as significant as
changing the amount of catalyst inside of the reactor.

\begin{figure}
\centering
\includegraphics{Ammonia2.png}
\caption{Heterogeneous Reactions Scale with Quantity of Catalyst, not
Reactor Volume!}
\end{figure}

Imagine we mix 1 mole of \(N_2\) and 3 moles of \(H_2\) in a 1L reactor
(\(C_{N_2} = 1M\) and \(C_{H_2} = 3M\)), and we add 1 gram of catalyst.
This catalyst achieves a constant \textbf{\emph{intensive}} rate of
reaction of 1 mole of extent per gram of catlayst per minute. For
simplicity, we are assuming that the rate of reaction is independent of
species concentration so that the rate of reaction is constant
throughout the reaction progress. This is not always going to be true.
We allow this system to react for one minute, at which point all of the
\(N_2\) and \(H_2\) are consumed. At this point in time, we will have
produced 2 moles of \(N \! H_3\) in the gas phase (\(C_{NH_3} = 2M\)).

Now we want to scale this process. In the second illustration of the
above figure, we increase the volume of the reactor and keep the initial
gas phase concentrations of \(N_2\) and \(H_2\) constant. We also keep 1
gram of catalyst constant; that catalyst still achieves a rate of
reaction of 1 mole of extent per gram of catalyst per minute; and we
allow the system to again react for 1 minute. In this case, since the
reaction only occurs on the catalyst, and we have kept the quantity of
catalyst the same (1 gram), we find that the total volume of the system
is irrelevant to determining the total quantity of \(N_2\) and \(H_2\)
that are converted into ammonia. After 1 minute elapses, we still
consume 1 mole of \(N_2\) and 3 moles of \(H_2\), and we still produce 2
moles of \(N \! H_3\). We won't increase the extensive quantity of
ammonia produced (the scale of the process) until we increase the amount
of catalyst, which is shown in the 3rd illustration. Only when we
increase the quantity of catalyst to 10 grams would we convert 10 moles
of \(N_2\) into 20 moles of \(N \! H_3\) in one minute.

Because of this scaling, for a reaction catalyzed by a solid, we often
see quantities reported per mass of catalyst. In many cases, ``mass of
catalyst'' is a good description of the size of the reactor, so we use
mass of catalyst as a basis for defining some intensive properties of
the system (like a reaction rate). If you look back at the specific way
we defined reaction rates in this example, you'll see that I was careful
to include the units of ``moles of extent per gram of catalyst per
minute.'' That is an example of an intensive reaction rate normalized to
a unit mass of catalyst. We reach a decision about that intensive basis
only after considering a sensible description of the size of our system.

    \hypertarget{heterogeneous-reactions-at-2d-interfaces}{%
\subsubsection{Heterogeneous reactions at 2D
interfaces}\label{heterogeneous-reactions-at-2d-interfaces}}

We might imagine that a reaction only occurs at the interface between
two phases. A common example might be a reaction between two immiscible
liquids. Alternatively, chemical vapor deposition is a widely used
technology for synthesizing solid materials. A good example is the
chemical vapor deposition of silane to produce silicon. This is a core
technology in the manufacturer of electronic components.

\[\text{Si} H_4 \ (g) \leftrightharpoons \text{Si} \ (s) + 2H_2 \ (g)\]

Here, you have a reaction where the gas phase silane decomposes at the
silicon surface to deposit more solid silicon. So, the reaction is only
occuring at the solid surface. In this case, the volume of silane is
largely irrelevant to the size of the silicon wafer that we are growing.
A better way to view the size of this system would be in terms of the
surface area of the silcon wafer as the reaction is not occuring
anywhere except for the surface of the silicon wafer. Here, we might see
reaction rates reported per unit surface area, and we'd find that the
total extensive quantity of silane that we're consuming scales with
surface area, i.e., interfacial surface area is probably a good basis
for the size of this system. Note that we also might use interfacial
surface area for heterogeneously catalyzed reactions as the
\textbf{\emph{surface reactions}} of heterogeneous catalysis are clearly
occuring on 2D surfaces. Finally, areal rates and properties are
appropriate for systems that involve immiscible liquids reacting at an
interface.

    \hypertarget{heterogenous-reactions-at-exposed-surface-atoms-active-sites}{%
\subsubsection{Heterogenous reactions at exposed surface atoms (active
sites)}\label{heterogenous-reactions-at-exposed-surface-atoms-active-sites}}

Catalytic reactions usually don't happen uniformly throughout the mass
or volume of a catalyst you put into a reactor. In reality, the
catalytic surface has ``active sites'' where the reaction occurs. Often,
these are atoms of expensive metals like Pt, Pd, or Ru. It is far more
cost effective to disperse these metals as isolated nanoparticles on an
inert, high surface area carrier like \(SiO_2\) or \(Al_2O_3\) than it
is to throw chunks of Pt into the reactor (which have a very low
percentage of their atoms on the surface--most Pt atoms in a chunk of Pt
are inside of the chunk of Pt, which does us no good from a catalysis
standpoint. Only the surface atoms are useful.)

So, revising our picture of ammonia synthesis, it probably looks more
like this:

\begin{figure}
\centering
\includegraphics{Ammonia3.png}
\caption{Actually\ldots they scale with the number of active sites more
than the mass of catalyst!}
\end{figure}

In these cases, it isn't the mass of catalyst or even the surface area
of catalyst that dicatates the extensive performance of the system: it
is the number of exposed active sites (Fe atoms in this example, but it
could be Pt atoms, Pd atoms, protons, etc). Often, this is the view of a
system we take in fundamental catalysis, where we are interested in
studying the system on the basis of a single active site, and so we
frequently report rates this way (these are usually called turnover
frequencies).

    \hypertarget{material-balances-on-reacting-systems}{%
\subsection{Material Balances on Reacting
Systems}\label{material-balances-on-reacting-systems}}

As reactions occur, we know that:

\begin{enumerate}
\def\labelenumi{\arabic{enumi}.}
\tightlist
\item
  Total mass is conserved
\item
  The total number of each element is conserved
\end{enumerate}

However, the nature of that mass changes form as the atoms break apart
and recombine in new ways. For example, 1 \(N_2\) and 3 \(H_2\)'s become
2 \(N\!H_3\)'s. Our ability to mathematically describe the change in
quantities of each species is critical to our ability to design a
chemical reactor. It allows us to answer questions like:

\begin{enumerate}
\def\labelenumi{\arabic{enumi}.}
\tightlist
\item
  What size of reactor do I need to achieve a certain conversion?
\item
  What conversion will the reactor achieve for a certain mass of
  catalyst?
\item
  How much time will it take for the reactor to achieve a certain yield
  of product?
\end{enumerate}

All of these are \textbf{\emph{extensive}} questions; in other words,
the answers to these questions depend on the scale of the system we are
considering. This is in contrast to an \textbf{\emph{intensive}}
property, which is independent of the scale of the system. To answer
these questions, we write \textbf{\emph{material balances}}, which allow
us to account for changes in quantities of each species that occur in
chemical reactors. Because balanced reactions generally describe changes
in molar quantities of reactants and products, we typically write
material balances for reactors in molar units. This makes it
straightforward for us to use stoichiometry and other properties of
reactions that are specified on a molar basis. Material balances (mole
balances) on reacting systems are probably the most important concept
covered in a reactor design course. Almost all of our reactor design
problems concerned with sizing a process or determining how long a
process will take require us to write material balances.

\hypertarget{developing-a-material-balance-for-a-reacting-system}{%
\subsubsection{Developing a Material Balance for a Reacting
System}\label{developing-a-material-balance-for-a-reacting-system}}

To start, we consider a system of arbitrary shape and ``size.'' Most of
us immediately think of Volume when we consider the ``size'' of a
reactor. Volume will always be important in determining the physical
space occupied by a reactor as well as its cost. However, as discussed
above, volume may not always be the best (mathematical) choice of basis
to describe the ``size'' of our system when reactions occur at
interfaces and catalyst surfaces. Before we get into that discussion,
let's look at general aspects of a mole balance that apply for a system
of any size or shape.

\begin{figure}
\centering
\includegraphics{MB1.png}
\caption{Illustration of an Arbitrary Chemical Reactor}
\end{figure}

The squiggly amorphous shape is our reactor -- it has an arbitrary
``size'' and an arbitrary shape. As of now, it is no type of reactor in
particular, it just represents a place where a reaction is occuring. If
we're going to write a material balance on that space, we have to
consider several things. First, in this course, we usually write
material balances for species rather than total material balances. With
this in mind, we will develop this balance for an arbirary species,
\(j\). Once we've established that we're going to write a balance on the
molar quantity of species \(j\) in that squiggly element of arbitrary
size and shape, we have to consider a few things that contribute to the
balance. One should recognize that each of these things is an extensive
quantity because it scales with the size of the system:

\begin{enumerate}
\def\labelenumi{\arabic{enumi}.}
\tightlist
\item
  The accumulation of moles of species \(j\) in the reactor as a
  function of time, \(d\!N_j/dt\).
\item
  The ``feed'' molar flowrate flow of species \(j\) into the reactor,
  \(F_{j,f}\).
\item
  The effluent molar flowrate of species \(j\) exiting the reactor,
  \(F_j\).
\item
  The rate of generation or consumption of species \(j\) by chemical
  reaction, \(G_j\).
\end{enumerate}

If we can account for those 4 quantities, we fully describe the material
balance (mole balance) on species \(j\) in our reactor.

I am using the generic notation that \(N_j\) is the moles of species
\(j\) in units of \textbf{\emph{moles of \(j\)}}, \(F_j\) and
\(F_{j,f}\) are molar flowrates of species \(j\), with units of
\textbf{\emph{moles of \(j\) per time}}, and \(G_j\) is the extensive
rate of ``generation'' by chemical reaction in units of
\textbf{\emph{moles of \(j\) per time}}.

We apply the following concept to describe a material balance on a
species:

\textbf{\{Accumulation of \(j\)\} = \{Flow of \(j\) in\} - \{Flow of
\(j\) out\} + \{Generation of \(j\) by chemical reactions\}}

If we substitute the labelled quantities in the figure above, we arrive
at the general material balance equation for a chemical reactor:

\[\frac{dN_j}{dt} = F_{j,f} - F_j + G_j\]

This equation applies (without making assumptions) to any reactor, of
any type, of any size, of any shape, and for any number of reactions
occuring inside of the reactor. It is a general starting point for
writing a balance on \textbf{\emph{any}} reacting system. The units of
each term in this equation are \textbf{\emph{moles per time}}. There is
no dimensional normalization. A general material balance of this type is
therefore an \textbf{\emph{extensive}} descrption of quantities of
species \(j\) within and moving through a reactor, and it will change
with the size of the system or the type of reactor you are considering.

\hypertarget{analyzing-the-parts-of-a-material-balance-on-species-j}{%
\subsubsection{\texorpdfstring{Analyzing the parts of a material balance
on species
\(j\)}{Analyzing the parts of a material balance on species j}}\label{analyzing-the-parts-of-a-material-balance-on-species-j}}

\hypertarget{influent-molar-flowrate}{%
\paragraph{Influent Molar Flowrate}\label{influent-molar-flowrate}}

Consider the molar flowrate of species \(j\) entering the reactor--we
usually refer to this as the ``feed'' molar flowrate of \(j\), and we
give it the symbol, \(F_{j,f}\). It is an extensive quantity and has
units of moles per time. Usually, as engineers--either in industrial
practice or in a research lab--we have some control over (or at least
some knowledge about) what goes into the reactor. Flowrates can be
controlled using a number of devices, and we have numerous ways of
either controlling or measuring the composition of an influent stream.
Sometimes, we'll have species molar flowrates (\(F_j\)) given to us
directly. In other cases, we will be given information about the total
molar or volumetric flowrate into the reactor as well as the
concentration or composition of each species in that feed stream. For
this reason, we will often see component molar flowrates specified in
terms of composition (mole fractions) and total molar flowrate, for
example:

\[F_{j,f} = y_{j,f}F_{T,f}\]

Where \(y_{j,f}\) is the mole fraction of \(j\) in the feed, and
\(F_{T,f}\) is the total molar flowrate of the feed stream.
Alternatively, we may be given feed concentrations and a total
volumetric flowrate into the reactor. There are practical reasons that
we often have access to this type of information -- species
concentrations are relatively easy to control or measure, and we often
control total volumetric flowrates of a fluid phase (using a pump)
rather than controlling molar flowrates directly.

\[F_{j,f} = C_{j,f}Q_{f}\]

Here, \(C_{j,f}\) is the concentration of \(j\) in the feed stream, and
\(Q_f\) is the total volumetric flowrate going into the reactor.

\hypertarget{effluent-molar-flowrate}{%
\paragraph{Effluent Molar Flowrate}\label{effluent-molar-flowrate}}

Now we consider the molar flowrate of species \(j\) leaving the reactor,
\(F_j\). More often than not, we will be describing our system in terms
of concentration and volumetric flowrates or composition (mole
fractions) and total molar flowrates, so we'll frequently need to relate
exit molar flowrates to these quantities. This is done exactly as above,
just redefined so that it applies at the ``exit'' of the reactor:

\[F_{j} = y_{j}F_T\]

or

\[F_{j} = C_{j}Q\]

It is worth considering that we will frequently need to measure the
composition or concentration of species at the conclusion of the
reaction. There is a very simple reason for this: despite knowing the
influent composition to a reactor, the reactor is designed to change
that composition, and we do not know a priori the extent of that change.
As such, the composition exiting the reactor is usually unknown to us,
so we must measure it using any number of analytical tools available.
Some common techniques for quantifying the composition of a mixture
include Gas Chromatography, Liquid Chromatography, NMR, FTIR, and
UV-Vis, though the appropriate analytical tool may change depending on
the specific system considered. Once we have quantified concentrations
and a volumetric flowrate, we can express molar flowrates in terms of
these quantities.

\hypertarget{accumulation-of-species}{%
\paragraph{Accumulation of species}\label{accumulation-of-species}}

This quantity represents the ``accumulation'' of species \(j\) in the
system with respect to time. All this means is that there is a change in
the total number of moles of species \(j\) in the system as time
changes. We can envision ``accumulation'' in many ways. For example, if
we fill a tank with water, we can see clearly that the liquid level and
the number of moles of water are changing in that tank as a function of
time. This would be an example of accumulation of water in the system,
and we can see it pretty clearly from the increase in liquid level!

\begin{figure}
\centering
\includegraphics{tankfill.png}
\caption{What constitutes accumulation?}
\end{figure}

Clearly, as time goes on, the number of moles of water in the tank are
increasing because the flow in is greater than the flow out. This is
``accumulation'' of water in the tank.

\hypertarget{accumulation-without-flow}{%
\subparagraph{Accumulation without
flow}\label{accumulation-without-flow}}

A less obvious example of species accumulation is the type of
accumulation that occurs in a batch reactor. Batch reactors are closed
systems where there is no accumulation of mass; however, the identity of
species will change with time. This is perhaps a non-obvious example of
accumulation. Consider the reaction:

\[A \ (l) \longrightarrow B \ (l)\]

Let's say we initially have a tank that will fill with 10 moles of
\(A\). It occupies a volume of 1L liter. The density of \(B\) is equal
to the density of \(A\), so as we convert \(A\) into \(B\), the system
volume does not change. Despite this, we have an increase in the number
of moles of \(B\) in that volume and a decrease in the number of moles
of \(A\) in that volume as time elapses.

In a batch reactor, total mass is constant. If density is also constant,
then the reactor volume must be constant.

\begin{figure}
\centering
\includegraphics{batchaccumulation.png}
\caption{If there is a reaction occuring in a batch reactor, there is
always accumulation of species!}
\end{figure}

We can see clearly that the number of moles of \(A\) and \(B\) are
changing in the tank as a function of time. This means there is
``accumulation'' of each species despite the fact that there is no flow
in or flow out, and the liquid level is constant. All accumulation means
in this context is that there is a change in the number of moles in the
system as time goes on.

Accumulation can be negative!

\hypertarget{the-rate-of-generation-or-consumption-by-chemical-reaction}{%
\paragraph{The Rate of Generation or Consumption by chemical
reaction}\label{the-rate-of-generation-or-consumption-by-chemical-reaction}}

The final term in the balance is probably the most nuanced. It is the
\textbf{\emph{extensive}} rate of generation (or consumption) of species
\(j\) by chemical reaction. Following the notation in Scott Fogler's
\emph{Elements of Chemical Reaction Engineering}, we will give it the
symbol \(G_j\), and it represents how quickly species \(j\) is either
consumed or produced by chemical reaction(s) that are occuring inside of
our reactor. It is definitely related to the rate of reaction, but
\textbf{\emph{it is not the same thing as the rate of reaction}}:

\[G_j \neq r_i\]

\(G_j\) is a \textbf{\emph{species specific term}}. It is an
\textbf{\emph{extensive}} quantity that represents the amount of species
\(j\) that is produced or consumed by chemical reaction in extensive
units of moles-per-time. It will change depending on the size of the
system and the particulars of the reactor you are using. In contrast,
the rate of reaction, \(r_i\), is a \textbf{\emph{reaction specific
term}}. It represents the rate of reaction \(i\) that is occuring at the
conditions inside of the reactor you are considering. Importantly, in
Chemical Reaction Engineering, the rate of a chemical reaction, \(r_i\),
is always an \textbf{\emph{intensive}} quantity that is scale and/or
system independent. In terms of what controls the intensive rate of
reaction, \(r_i\):

\[r_i = f(T, P, \chi_j)\]

To expand that concept: once we decide on what reaction we're
considering and what (if any) catalyst it uses, \textbf{\emph{the rate
of that reaction, \(r_i\), is entirely determined by the temperature,
pressure, and composition of the reacting mixture}}. The rate of
reaction, \(r_i\), does not depend on stoichiometry, nor does it depend
on the type, size, or shape of a reactor. It might surprise you, but the
rate of reaction also has no fundamental dependence on important aspects
of reactor design, like mixing. The rate of reaction is always entirely
determined by temperature, pressure, and composition. In contrast,
stoichiometry, shape, size, type of reactor, extent of mixing, etc. will
all influence the extensive rate of generation by chemical reaction,
\(G_j\):

\[G_j = f(T, P, \chi_j, \text{reactor type, reactor size, mixing, etc.})\]

These two quantities are related, but they are not the same thing. We'll
formally develop a generation rate in terms of reaction rates in
\href{https://jqbond.github.io/CEN587/Notebooks/587-N14.html}{Notebook
14}.


    % Add a bibliography block to the postdoc
    
    
    
\end{document}
